% 编译方式: xelatex SL_fractional_left_definite.tex
\documentclass[UTF8]{ctexart}
\usepackage{amsmath, amssymb, amsthm}
\usepackage[margin=2.5cm]{geometry}
\usepackage[hyphens]{url}
\usepackage[hidelinks]{hyperref}
\usepackage{booktabs}
\emergencystretch 3em
\allowdisplaybreaks

\newtheorem{theorem}{定理}[section]
\newtheorem{proposition}[theorem]{命题}
\newtheorem{lemma}[theorem]{引理}
\newtheorem{corollary}[theorem]{推论}
\newtheorem{remark}[theorem]{注}
\newtheorem{definition}[theorem]{定义}

\title{稀疏多项式族的完整非负阶稠密窗口与成员阈值}
\author{研究证明文档 (项目: BVE research), 第六轮作者修订}
\date{2026-08-05 初稿; 2026-09-20 勘误; 2026-09-21 补证}

\begin{document}
\maketitle

\begin{abstract}
固定 $c>0$, 对移位 Krein Laplacian 的真正左定幂域
$H_c^s=D(K_c^{s/2})$, 本文证明原稀疏多项式族在全部
$0\le s<7/2$ 中稠密. 每个非仿射命名成员的非负阶成员阈值恰为
$s<7/2$, 而 $1,x$ 属于全部非负幂域. 证明逐一覆盖偶, 奇两支的所有指标,
给出归一化谱系数的非零边界首项及端点发散, 再用行列式为 $15360$ 的
四迹修正构造 $D(K_c^2)$ 中的多项式图核心, 经谱截断向下传输.
这补齐旧稿未判定的 $3<s<7/2$, 不恢复已经撤回的全阶命名族断言.
本版是对所供 R6-C01 候选的作者展开和核查, 不是独立验收回执或完整 Lean 形式化.
\end{abstract}

\tableofcontents

\section{对象, 问题合同与来源}
全部空间取复数域, 内积 $\langle f,g\rangle_2=\int_{-1}^1f\overline g$
对第一变量线性, 线性包只含有限线性组合. 固定 $c>0$, 记 $I=(-1,1)$,
$\Delta f=f(1)-f(-1)$, 并定义
\begin{equation}\label{eq:operator-domain}
 \begin{gathered}
 K_cf=-f''+cf,\qquad D(K_c)=\{f\in H^2_{\rm Sob}(I):\mathcal Bf=0\},\\
 \mathcal Bf=\binom{f'(1)-\Delta f/2}{f'(-1)-\Delta f/2}.
 \end{gathered}
\end{equation}
$H^r_{\rm Sob}$ 表示普通 Sobolev 空间. 左定空间记为
$H_c^s=D(K_c^{s/2})$ ($s\ge0$), 范数为
$\|f\|_{s,c}=\|K_c^{s/2}f\|_2$. 以下固定 $c$ 后也简写为 $H^s,\|f\|_s$.

原族的指标集是 $\mathcal D=\{0,1\}\cup\{4,5,6,\ldots\}$, 其中
\begin{equation}\label{eq:original-family}
 \begin{gathered}
 p_0=1,\quad p_1=x,\\
 p_{2m}=x^{2m}-\frac m{m-1}x^{2m-2},\qquad
 p_{2m+1}=x^{2m+1}-\frac m{m-1}x^{2m-1},\quad m\ge2.
 \end{gathered}
\end{equation}
本文只证明这个\emph{完整原族}的非负阶成员性与稠密性. 各固定多项式的估计
常数可以依赖它的次数, 嵌入常数可以依赖 $c$; 不断言 $c\downarrow0$ 时一致.

修订基线为 \texttt{92d46e99e36cc0a74e8f7bb13430749c222ee54f}.
所供 \cite{r6candidate} 是候选论证来源. 算子基础实际核对了
\cite[第2节, R5-OP]{r5}, 代数消元实际核对了 \cite[第2.1节]{h2}
及 \cite{dc}; 多项式逼近与迹修正还参照了 \cite[STEP7--STEP8]{a7} 的证明.
下文重新给出本题所需的算子, 谱系和四阶图核心论证, 不以旧审计标签为前提,
也不使用已经被反证的一般投影稀疏族命题. R5 原文及其他历史证明保持原字节.
本文件的本轮修订前 TeX/PDF 保存在
\path{research/artifacts/proof-audit-round6-20260921/before/docs/}.

\section{算子基础, 完整谱系与标度}
\subsection{正性, 满射, 自伴性与紧逆}
\begin{proposition}[算子基础]\label{prop:operator}
\eqref{eq:operator-domain} 定义稠密自伴算子, $K_c\ge cI$,
且 $K_c^{-1}:L^2\to L^2$ 紧. 在 $D(K_c)$ 上有
\begin{equation}\label{eq:h2-estimate}
 \|f\|_{H^2_{\rm Sob}}\le C_c\|K_cf\|_2.
\end{equation}
\end{proposition}
\begin{proof}
由分部积分和 $\Delta f=\int_{-1}^1f'$,
\[
 \langle K_cf,f\rangle_2
 =\|f'\|_2^2-\frac12|\Delta f|^2+c\|f\|_2^2
 \ge c\|f\|_2^2.
\]
因此 $\|f\|_2\le c^{-1}\|K_cf\|_2$,
$\|f''\|_2\le2\|K_cf\|_2$. 为补足一阶导数, 写
$f(x)=A+Bx+r(x)$, 其中 $r(x)=\int_0^x(x-t)f''(t)\,dt$.
两半区间上的 Cauchy--Schwarz 不等式给出
\[
 \|r\|_2\le\frac1{\sqrt{12}}\|f''\|_2,\quad
 \|r'\|_2\le\frac1{\sqrt2}\|f''\|_2,\quad
 \|f'\|_2\le\sqrt3\|f\|_2+
 \left(\frac12+\frac1{\sqrt2}\right)\|f''\|_2.
\]
最后一式用到了 $\|A+Bx\|_2^2=2|A|^2+(2/3)|B|^2$.
这证明 \eqref{eq:h2-estimate}.

为直接证明满射, 令 $a=\sqrt c$, $d_a=a\cosh a-\sinh a$.
因 $d_0=0$, $d_a'=a\sinh a>0$, 有 $d_a>0$. 对 $h\in L^2$, 置
\[
 v_h(x)=-\frac1a\int_0^x\sinh(a(x-t))h(t)\,dt.
\]
负 $x$ 时使用有向积分. 积分核与其一阶导数有界,
$v_h''=cv_h-h\in L^2$, 故 $v_h\in H^2_{\rm Sob}$ 且 $K_cv_h=h$.
两个齐次解的边界矩阵为
\[
 \begin{pmatrix}\mathcal B\cosh(ax)&\mathcal B\sinh(ax)\end{pmatrix}
 =\begin{pmatrix}a\sinh a&d_a\\-a\sinh a&d_a\end{pmatrix},
 \qquad\det=2a\sinh a\,d_a>0.
\]
所以唯一的线性组合修正 $v_h$ 的边界残差, 给出 $f\in D(K_c)$,
$K_cf=h$. 单射性由正性得到.

$C_c^\infty(I)\subset D(K_c)$ 给出稠密定义性;
边界形式 $[-f'\overline g+f\overline{g'}]_{-1}^1$ 对域内 $f,g$ 为零,
所以算子对称. 若 $u\in D(K_c^*)$, 取 $v=K_c^{-1}K_c^*u\in D(K_c)$.
对所有 $f\in D(K_c)$,
$\langle K_cf,u-v\rangle_2=\langle f,K_c^*u-K_cv\rangle_2=0$.
满射性迫使 $u=v$, 故 $K_c^*=K_c$.
由 \eqref{eq:h2-estimate}, $K_c^{-1}$ 将 $L^2$ 有界集送入 $H^2_{\rm Sob}$ 有界集.
后一有界集的函数一致有界, 且
$|f(x)-f(y)|\le\|f'\|_2|x-y|^{1/2}$.
Arzelà--Ascoli 定理给出一致收敛子列, 因而也在 $L^2$ 收敛, 证明紧逆.
这里用到的端点及一致界均可由
\begin{equation}\label{eq:trace-bound}
 |u(t)|\le2^{-1/2}\|u\|_2+\sqrt2\|u'\|_2,\qquad
 u\in H^1_{\rm Sob}(I),\quad -1\le t\le1
\end{equation}
得到: 将 $u(t)$ 与平均值比较, 再积分 $u'$ 即可.
\end{proof}

\subsection{全部本征函数及归一化}
紧正自伴逆的谱定理给出 $L^2$ 的完备正交本征系; 其零核为 $\{0\}$,
所以没有遗漏的正交补. 下面解出所有本征函数, 而不是只列出若干本征对.
正性排除 $\lambda<c$. 当 $\lambda=c$, 方程 $-f''=0$ 给出全部仿射函数,
均满足边界条件, 故重数恰为二. 当 $\lambda>c$, 令 $\omega=\sqrt{\lambda-c}$.
反射 $f(x)\mapsto f(-x)$ 保持算子域并与 $K_c$ 交换, 可分偶, 奇求解.
偶部在 $[0,1]$ 满足 $f'(0)=f'(1)=0$, 故为 $\cos(k\pi x)$, $k\ge1$.
奇部满足 $f(0)=0$, $f'(1)=f(1)$, 故为 $\sin(\mu x)$,
$\mu\cos\mu=\sin\mu$. 极点不满足这个等式.
在 $(k\pi,(k+1/2)\pi)$ 上, $\tan\mu-\mu$ 从 $-k\pi$ 严格增加到 $+\infty$,
导数是 $\tan^2\mu>0$; 对 $k\ge1$ 恰有一根 $\mu_k$.
在 $(0,\pi/2)$ 上它严格为正, 其余正半周期 $\tan\mu<0$,
所以没有其他正根. 零频奇模态已经计入 $x$.
于是完整谱表为
\[
\begin{array}{c|c|c}
 \phi & \lambda & \|\phi\|_2^2\\ \hline
 1 & c & 2\\
 x & c & 2/3\\
 \cos(k\pi x),\ k\ge1 & (k\pi)^2+c & 1\\
 \sin(\mu_kx),\ k\ge1 & \mu_k^2+c & \mu_k^2/(1+\mu_k^2)
\end{array}
\]
不同特征值的正交性由自伴性给出, $1,x$ 由奇偶性正交.
直接积分还核对了容易遗漏的低频奇模态和正弦范数:
\[
 \int_{-1}^1x\sin(\mu_kx)\,dx
 =\frac{2(\sin\mu_k-\mu_k\cos\mu_k)}{\mu_k^2}=0,\qquad
 \|\sin(\mu_kx)\|_2^2=1-\frac{\sin(2\mu_k)}{2\mu_k}
 =\frac{\mu_k^2}{1+\mu_k^2}.
\]
明确的归一化为
\begin{equation}\label{eq:normalized-modes}
 \begin{gathered}
 e_0=\frac1{\sqrt2},\quad e_1=\sqrt{\frac32}x,\quad
 e_k^{\rm e}=\cos(k\pi x),\\
 e_k^{\rm o}=\frac{\sqrt{1+\mu_k^2}}{\mu_k}\sin(\mu_kx),\quad k\ge1.
 \end{gathered}
\end{equation}
记 $N_k=\mu_k/\sqrt{1+\mu_k^2}$, 则
$\|\sin(\mu_kx)\|_2=N_k\ge N_1>0$,
$\sin\mu_k=(-1)^kN_k$, $N_k\to1$.
这组谱系与 \cite[第2节, 式(7)--(9)]{axioms} 的模型一致;
本节已从定义域与紧逆独立推出完备性.

\subsection{非负幂域和负阶完备化的区别}
把 \eqref{eq:normalized-modes} 合并枚举为 $\{e_j\}$, 对 $s\ge0$,
谱演算给出真正的幂域判据
\begin{equation}\label{eq:spectral-norm}
 f\in H^s\quad\Longleftrightarrow\quad
 \sum_j\lambda_j^s|\langle f,e_j\rangle|^2<\infty,\qquad
 \|f\|_{H^s}^2=\sum_j\lambda_j^s|\langle f,e_j\rangle|^2.
\end{equation}
若使用未归一化的 $\phi_j$, 每项必须除以 $\|\phi_j\|_2^2$.

对 $s<0$, 定义 $H^s$ 为 $L^2$ 在范数 $\|K_c^{s/2}f\|_2$ 下的完备化,
即加权序列空间 $\{(a_j):\sum_j\lambda_j^s|a_j|^2<\infty\}$.
虽然有界算子 $K_c^{s/2}$ 在 $L^2$ 上处处有定义, 这个较弱范数不与
$L^2$ 范数等价: $\|\cos(n\pi x)\|_s=((n\pi)^2+c)^{s/2}\to0$,
而其 $L^2$ 范数恒为一. 事实上对任意 $s<0$, 取余弦模态系数
$a_n=n^{-1/2}$, 则 $\sum\lambda_n^s/n<\infty$ 而 $\sum1/n=\infty$;
其有限和在弱范数下 Cauchy, 在 $L^2$ 中没有该范数的极限.
这同时说明必须作完备化. 在这个完整标度上, 系数映射
\[
 \mathcal K_c:(a_j)\longmapsto(\lambda_ja_j),\qquad
 \mathcal K_c:H^s\longrightarrow H^{s-2}
\]
对每个实数 $s$ 是等距满射, 逆映射逐项乘以 $\lambda_j^{-1}$.
它在 $D(K_c)$ 上与原算子相同, 在其它阶上表示标度延拓, 不能默认
对所有元素都按经典二阶导数作用. 保留这一负阶定义, 但本轮稠密性定理
只陈述非负阶; 以下证明不使用负阶空间.

\section{全部命名成员的精确阈值}
\begin{lemma}[所有边界相容多项式的系数上界]\label{lem:coefficient-upper}
对任意固定多项式 $g$, 两支高频谱系数均为 $O_g(k^{-2})$.
若 $p\in\mathbb C[x]\cap D(K_c)$, 两支归一化系数均为 $O_{p,c}(k^{-4})$,
从而 $p\in H_c^s$ 对所有 $0\le s<7/2$ 成立.
\end{lemma}
\begin{proof}
记 $C_k(g)=\langle g,\cos(k\pi x)\rangle_2$,
$S_k(g)=\langle g,\sin(\mu_kx)\rangle_2$, $a_k=k\pi$.
分部积分两次, 得到精确恒等式
\begin{align}
 C_k(g)&=\frac{(-1)^k\bigl(g'(1)-g'(-1)\bigr)}{a_k^2}
          -\frac{C_k(g'')}{a_k^2},\label{eq:cos-ibp}\\
 S_k(g)&=-\frac{\Delta g\cos\mu_k}{\mu_k}
 +\frac{(g'(1)+g'(-1))\sin\mu_k}{\mu_k^2}
 -\frac{S_k(g'')}{\mu_k^2}\notag\\
 &=\frac{\sin\mu_k}{\mu_k^2}
 \bigl(g'(1)+g'(-1)-\Delta g\bigr)-\frac{S_k(g'')}{\mu_k^2}.
 \label{eq:sin-ibp}
\end{align}
末项积分的绝对值不超过 $\|g''\|_{L^1}$, 且
$k\pi<\mu_k<(k+1/2)\pi$, 所以两式均为 $O_g(k^{-2})$.
这些等式使用的是 $L^2$ 配对, 不要求 $g\in D(K_c)$.

若 $p\in D(K_c)$ 是多项式, 则 $K_cp$ 也是多项式.
对任一高频本征函数 $\phi$ 及其特征值 $\lambda$,
自伴性给出
\[
 \langle p,\phi\rangle_2=\lambda^{-1}\langle K_cp,\phi\rangle_2.
\]
由 $\lambda\asymp_c k^2$ 及前面的上界, 未归一化系数为 $O_{p,c}(k^{-4})$.
除以 $N_k\ge N_1>0$ 仍保持该阶数. 低频只有两项,
高频的谱范数平方受常数乘 $\sum_{k\ge1}k^{2s-8}$ 控制,
故在 $s<7/2$ 收敛. 此处不要求 $K_cp\in D(K_c)$.
\end{proof}

\begin{theorem}[逐一成员阈值]\label{thm:all-thresholds}
对每个固定 $c>0$ 和每个 $n\in\mathcal D\setminus\{0,1\}$,
\[
 p_n\in H_c^s\quad\Longleftrightarrow\quad 0\le s<7/2
 \qquad(s\ge0).
\]
$p_0,p_1$ 属于全部非负幂域. 更精确地, 对每个整数 $m\ge2$,
\begin{align}
 \langle p_{2m},e_k^{\rm e}\rangle_2
 &=-\frac{8m(4m-5)(-1)^k}{(k\pi)^4}+O_m(k^{-6}),\label{eq:even-leading}\\
 \langle p_{2m+1},e_k^{\rm o}\rangle_2
 &=-\frac{8m(4m-3)(-1)^k}{\mu_k^4}+O_m(k^{-6}).\label{eq:odd-leading}
\end{align}
相反奇偶的高频系数恒为零.
\end{theorem}
\begin{proof}
直接求导得到
\[
 p_{2m}'(\pm1)=0,\quad\Delta p_{2m}=0,\qquad
 p_{2m+1}'(\pm1)=p_{2m+1}(1)=-\frac1{m-1}
 =\frac{\Delta p_{2m+1}}2.
\]
所以全部原成员属于 $D(K_c)$, 引理 \ref{lem:coefficient-upper} 已给出窗口内成员性.
现在对所有 $m\ge2$ 计算不能相消的下一层边界数据:
\begin{align}
 p_{2m}'''(1)
 &=2m(2m-1)(2m-2)
 -\frac m{m-1}(2m-2)(2m-3)(2m-4)\notag\\
 &=2m\bigl((2m-1)(2m-2)-(2m-3)(2m-4)\bigr)\notag\\
 &=4m(4m-5)>0,\label{eq:even-boundary}\\
 p_{2m+1}'''(1)-p_{2m+1}''(1)
 &=(2m+1)(2m)(2m-2)
 -\frac m{m-1}(2m-1)(2m-2)(2m-4)\notag\\
 &=2m\bigl((2m+1)(2m-2)-(2m-1)(2m-4)\bigr)\notag\\
 &=4m(4m-3)>0.\label{eq:odd-boundary}
\end{align}
这些是无次数上界的代数恒等式; $m=2$ 时相应低次单项式的高阶导数为零,
上述乘积中的零因子正好覆盖此情况.

由于第一层边界项为零, \eqref{eq:cos-ibp} 连用两次给出
\[
 C_k(p_{2m})
 =-\frac{2(-1)^kp_{2m}'''(1)}{a_k^4}
   +\frac{C_k(p_{2m}^{(4)})}{a_k^4}.
\]
对末项再用引理中的任意多项式估计, 得 $O_m(k^{-6})$,
于是得到 \eqref{eq:even-leading}.
同理, $p_{2m+1}''$ 为奇函数, 故
\[
 (p_{2m+1}'')'(1)+(p_{2m+1}'')'(-1)-\Delta p_{2m+1}''
 =2\bigl(p_{2m+1}'''(1)-p_{2m+1}''(1)\bigr).
\]
由 \eqref{eq:sin-ibp} 连用两次,
\[
 S_k(p_{2m+1})
 =-\frac{8m(4m-3)\sin\mu_k}{\mu_k^4}
   +\frac{S_k(p_{2m+1}^{(4)})}{\mu_k^4}.
\]
除以 $N_k$, 使用 $\sin\mu_k/N_k=(-1)^k$ 及 $N_k\ge N_1>0$,
便得 \eqref{eq:odd-leading}. 这明确处理了正弦归一化, 并非从 $p_4$ 外推.

每个固定 $m$ 的首项常数均严格非零, 余项比首项小 $O_m(k^{-2})$.
因此存在 $b_m>0,k_m$, 使相应归一化系数的绝对值对 $k\ge k_m$
至少为 $b_m k^{-4}$. 结合特征值与 $k^2$ 可比, 谱范数平方尾部至少为
正的常数乘 $\sum_{k\ge k_m}k^{2s-8}$.
在 $s=7/2$ 这是调和级数, 在 $s>7/2$ 也发散.
最后 $K_c(a+bx)=c(a+bx)$, 故仿射函数的谱只在特征值 $c$ 上,
属于全部非负幂域.
\end{proof}

最低奇例也有精确式
$\langle p_5,\sin(\mu_kx)\rangle_2=-80\sin\mu_k/\mu_k^4$,
因为 $p_5^{(4)}=120x$ 与所有正频奇模态正交.
这只是全指标证明的交叉核对, 不是替代证明.

\section{四阶相容多项式图核心}
\subsection{原族的代数包与第二层算子域}
\begin{lemma}[精确的代数线性包]\label{lem:algebra}
令 $\mathcal P_K=\operatorname{span}_{\mathbb C}\{p_n:n\in\mathcal D\}$, 则
\[
 \mathcal P_K=\mathbb C[x]\cap D(K_c).
\]
\end{lemma}
\begin{proof}
每个生成元满足边界条件, 给出一侧包含.
反过来, 对右侧多项式按最高次数 $n\ge4$ 减去首项系数乘 $p_n$,
边界条件保留且次数严格下降. 最后余式为 $A+Bx+Dx^2+Ex^3$.
因为
\[
 \mathcal B1=\mathcal Bx=0,\quad
 \mathcal Bx^2=\binom2{-2},\quad\mathcal Bx^3=\binom22,
\]
两个残差向量线性无关, 迫使 $D=E=0$, 完成有限消元.
这里没有使用 $\mathcal P_K=\mathbb C[x]$ 这一错误等式.
\end{proof}

\begin{lemma}[真正的平方域]\label{lem:square-domain}
\begin{equation}\label{eq:square-domain}
 \begin{gathered}
 D(K_c^2)=\{f\in H^4_{\rm Sob}(I):\mathcal Bf=0,\ \mathcal B(f'')=0\},\\
 K_c^2f=f^{(4)}-2cf''+c^2f.
 \end{gathered}
\end{equation}
\end{lemma}
\begin{proof}
按乘积域定义, $f\in D(K_c^2)$ 当且仅当 $f,K_cf\in D(K_c)$.
于是 $f''=cf-K_cf\in H^2_{\rm Sob}$, 故 $f\in H^4_{\rm Sob}$.
由 $\mathcal B(K_cf)=c\mathcal Bf-\mathcal B(f'')$ 得两个迹条件.
反向取右侧的 $f$, 则 $f\in D(K_c)$, $K_cf\in H^2_{\rm Sob}$
且 $\mathcal B(K_cf)=0$, 故 $f\in D(K_c^2)$.
只有在这一步确认乘积域后, 才使用显示的四阶微分表达式.
\end{proof}

\subsection{四迹修正的有界右逆}
定义
\[
 T:H^4_{\rm Sob}(I)\longrightarrow\mathbb C^4,\qquad
 Tf=(\mathcal Bf,\mathcal B(f''))^{\mathsf T},
\]
排列顺序为先 $f$ 的 $+$, $-$ 残差, 再 $f''$ 的 $+$, $-$ 残差.
将 \eqref{eq:trace-bound} 用于 $f,f',f'',f'''$,
得到 $\|Tf\|_{\mathbb C^4}\le C_T\|f\|_{H^4_{\rm Sob}}$.
直接对 $x^2,x^3,x^4,x^5$ 求导, 矩阵为
\begin{equation}\label{eq:four-trace-matrix}
 B=\begin{pmatrix}
 2&2&4&4\\
 -2&2&-4&4\\
 0&0&24&40\\
 0&0&-24&40
 \end{pmatrix},\qquad
 \det B=(8)(1920)=15360\ne0.
\end{equation}
对 $z\in\mathbb C^4$, 定义多项式右逆
\[
 Rz=(x^2,x^3,x^4,x^5)B^{-1}z.
\]
于是 $TR=I_{\mathbb C^4}$, 且
\[
 \begin{gathered}
 \|Rz\|_{H^4_{\rm Sob}}\le C_R\|z\|_{\mathbb C^4},\\
 C_R=\left(\sum_{\ell=2}^5\|x^\ell\|_{H^4_{\rm Sob}}^2\right)^{1/2}
          \|B^{-1}\|_{\mathbb C^4\to\mathbb C^4}<\infty.
 \end{gathered}
\]
因此修正是固定的有界有限秩操作, 与待逼近函数及逼近序号无关.

\subsection{普通四阶 Sobolev 逼近及图范数控制}
\begin{lemma}[四次积分的多项式逼近]\label{lem:h4-approximation}
多项式在 $H^4_{\rm Sob}(I)$ 中稠密.
\end{lemma}
\begin{proof}
给定 $f\in H^4_{\rm Sob}$, 由连续函数在 $L^2$ 中稠密和 Weierstrass 定理,
取复多项式 $h_j\to f^{(4)}$ 于 $L^2$. 令
\[
 r_j(x)=\sum_{\ell=0}^3\frac{f^{(\ell)}(-1)}{\ell!}(x+1)^\ell
       +\int_{-1}^x\frac{(x-t)^3}{6}h_j(t)\,dt.
\]
这是多项式. $f^{(\ell)}$ ($0\le\ell\le3$) 均有绝对连续代表,
所以反复用微积分基本定理, $f$ 有同样的式子, 将 $h_j$ 换成 $f^{(4)}$ 即可.
设 $w_j=r_j-f$, 则 $w_j^{(4)}=h_j-f^{(4)}$, 低阶左端点迹均为零.
对 $0\le r\le3$ 的积分核使用 Cauchy--Schwarz, 再积分三角区域, 得
\[
 \|w_j^{(r)}\|_2\le
 \frac{2^{4-r}}{(3-r)!\sqrt{(7-2r)(8-2r)}}\,\|h_j-f^{(4)}\|_2.
\]
因为平方核积分恰为
$2^{8-2r}/((3-r)!^2(7-2r)(8-2r))$.
所有四个低阶导数及第四阶导数都收敛, 即 $r_j\to f$ 于 $H^4_{\rm Sob}$.
\end{proof}

\begin{proposition}[四阶多项式图核心]\label{prop:h4-core}
令 $\mathcal C_4=\mathbb C[x]\cap D(K_c^2)$. 则
\[
 \mathcal C_4\subset\mathcal P_K,\qquad
 \overline{\mathcal C_4}^{\,H_c^4}=H_c^4.
\]
它也在 $K_c^2$ 的通常图范数
$(\|f\|_2^2+\|K_c^2f\|_2^2)^{1/2}$ 中为核心.
\end{proposition}
\begin{proof}
包含关系由 $D(K_c^2)\subset D(K_c)$ 及引理 \ref{lem:algebra} 得到.
对 $f\in D(K_c^2)$, 取引理 \ref{lem:h4-approximation} 的 $r_j$, 置
\[
 q_j=r_j-RT r_j.
\]
因 $Tf=0$, $TR=I$, 有 $Tq_j=0$. 由 \eqref{eq:square-domain},
$q_j\in\mathcal C_4$, 并且
\[
 \|q_j-f\|_{H^4_{\rm Sob}}
 \le(1+C_RC_T)\|r_j-f\|_{H^4_{\rm Sob}}\longrightarrow0.
\]
在已确认的域内, \eqref{eq:square-domain} 给出
\begin{align*}
 \|K_c^2(q_j-f)\|_2
 &\le \|(q_j-f)^{(4)}\|_2+2c\|(q_j-f)''\|_2+c^2\|q_j-f\|_2\\
 &\le(1+2c+c^2)\|q_j-f\|_{H^4_{\rm Sob}}\longrightarrow0.
\end{align*}
这就是 $H_c^4$ 收敛. 此外谱下界给出
$\|u\|_2\le c^{-2}\|K_c^2u\|_2$, 故该范数与通常图范数等价.
整个论证只用了四阶域恒等式, 没有假定全部高阶或分数阶域的边界刻画.
\end{proof}

\begin{remark}[命名成员与相容组合的区别]\label{rem:cancellation}
尽管 $p_4,p_6\notin H_c^4$, 其有限线性组合可以消去第二层边界残差:
\[
 q=p_6-\frac72p_4=x^6-5x^4+7x^2,\qquad
 q'(1)=q'(-1)=q'''(1)=q'''(-1)=0.
\]
$q,q''$ 为偶函数, 所以两个值差也为零, 即 $Tq=0$,
$q\in\mathcal C_4$. 例如 $p_6'''(1)=84$, $(7/2)p_4'''(1)=84$,
相应 $k^{-4}$ 首项也正好抵消.
表达式 $q$ 是在多项式空间中作代数相加后再检验域,
不在 $H_c^4$ 中给域外的 $p_4,p_6$ 各自赋范数.
因此 $\operatorname{span}\{p_n:p_n\in H_c^4\}=\operatorname{span}\{1,x\}$
与 $\mathcal P_K\cap H_c^4=\mathcal C_4$ 是不同的空间;
前者为二维闭空间, 后者是无限维空间 $H_c^4$ 的稠密子空间.
\end{remark}

\section{完整非负阶窗口与谱截断传输}
\begin{theorem}[完整原族的非负阶稠密窗口]\label{thm:main}
对每个固定 $c>0$ 和 $0\le s<7/2$,
\[
 \{p_n:n\in\mathcal D\}\subset H_c^s,\qquad
 \overline{\operatorname{span}\{p_n:n\in\mathcal D\}}^{\,H_c^s}=H_c^s.
\]
对 $s\ge7/2$, 完整原族不是 $H_c^s$ 的子集.
此时逐个筛选已经属于该空间的命名成员, 只剩 $1,x$, 其线性包不稠密.
\end{theorem}
\begin{proof}
成员性及其障碍已经由定理 \ref{thm:all-thresholds} 证明.
固定 $s\in[0,7/2)$, $f\in H_c^s$ 和 $\varepsilon>0$.
令 $E$ 为 $K_c$ 的谱测度, $f_N=E([c,N])f$ ($N\ge c$).
因谱支集有界, $f_N\in H_c^4$; 且
\[
 \|f-f_N\|_{s,c}^2
 =\int_{(N,\infty)}\lambda^s\,d\|E_\lambda f\|^2\longrightarrow0.
\]
对任意 $u\in H_c^4$, 谱下界给出固定 $c$ 的连续嵌入
\begin{equation}\label{eq:h4-embedding}
 \|u\|_{s,c}^2
 =\int_{[c,\infty)}\lambda^s\,d\|E_\lambda u\|^2
 \le c^{s-4}\|u\|_{4,c}^2,\qquad 0\le s\le4.
\end{equation}
先选 $N$ 使截断误差小于 $\varepsilon/2$.
再由命题 \ref{prop:h4-core} 取 $q\in\mathcal C_4\subset\mathcal P_K$, 使
\[
 \|f_N-q\|_{4,c}<\frac{\varepsilon}{2c^{(s-4)/2}}.
\]
则 $\|f-q\|_{s,c}<\varepsilon$. 由于每个命名成员确实属于当前 $H_c^s$,
这就是目标空间内部完整原族的稠密性.
当 $s\ge7/2$, 仅 $1,x$ 单独合格, 它们张成的二维空间是闭的;
$H_c^s$ 含全部本征函数, 为无限维, 所以此筛选族不稠密.
\end{proof}

\begin{remark}[为何补齐窗口而不恢复全阶推广]
核心的稠密性传输本身还给出
$\overline{\mathcal C_4}^{\,H_c^s}=H_c^s$ ($0\le s\le4$).
但在 $s\ge7/2$ 时, $\mathcal P_K$ 已不是 $H_c^s$ 内的线性子空间,
不能把域交稠密性说成全部命名成员在该空间中构成稠密族.
新结论同时需要全成员的 $s<7/2$ 上界和四阶图核心, 不是从
$H^3$ 稠密性向更强的范数外推.

旧稿由 \cite{h3} 的 $H^1$ 矩跳跃主链所得 $H_c^3$ 稠密性仍提供独立的低阶路线.
其中 $D(K_c^{1/2})=H^1_{\rm Sob}$ 是该模型的形式域,
$K_cp_n$ 为多项式, 所以 $p_n\in H_c^3$.
对 $0\le s\le3$, 旧嵌入
$\|u\|_{s,c}\le c^{(s-3)/2}\|u\|_{3,c}$
配合先截断再逼近, 给出原来的低阶结论.
当 $s>3$, 该方向的有界嵌入不再成立: 高频本征函数的范数比
$\|e_j\|_{s,c}/\|e_j\|_{3,c}=\lambda_j^{(s-3)/2}\to\infty$.
本轮上移至相容的四阶逼近, 才补上 $3<s<7/2$.

本文不由此扩展到任意受约束子空间, 任意删项族或稳定数值基;
$H_c^4$ 以上的其他域交结论和负阶稠密性也不在本次定理的量词内.
\end{remark}

\section{反例与撤回的证明机制}
\begin{proposition}[精确成员阈值]\label{prop:thresholds}
对每个 $c>0$, $t,s\ge0$,
\[
 x^2\in H^t\iff t<3/2,\qquad p_4\in H^s\iff s<7/2.
\]
因此 $s\ge7/2$ 时原稀疏族甚至不是 $H^s$ 的子集.
\end{proposition}
\begin{proof}
对 $a=n\pi$, 两次分部积分给出
\[
 \int_{-1}^1x^2\cos(ax)\,dx=\frac{4(-1)^n}{a^2},\qquad
 \int_{-1}^1x^4\cos(ax)\,dx
 =\frac{8(-1)^n}{a^2}-\frac{48(-1)^n}{a^4}.
\]
所以 $p_4$ 的 $a^{-2}$ 项恰好相消. 与未归一化常数模态的内积分别为 $2/3$ 与
$-14/15$, 所有奇模态的投影为零. 允许级数值为 $+\infty$, 得到
\begin{align*}
 \|x^2\|_t^2&=\frac{2c^t}{9}
 +16\sum_{n\ge1}\frac{((n\pi)^2+c)^t}{(n\pi)^4},\\
 \|p_4\|_s^2&=\frac{98c^s}{225}
 +2304\sum_{n\ge1}\frac{((n\pi)^2+c)^s}{(n\pi)^8}.
\end{align*}
尾项分别与 $n^{2t-4}$ 和 $n^{2s-8}$ 等价. 正项级数判别给出上述
严格阈值, 两个端点均为调和级数型发散. 这里 $+\infty$ 只表示不属于
相应空间, 不是空间元素的有限范数. \qed
\end{proof}

保留第一轮的独立算子域反例作为一致性检查:
$p_4\in D(K_c)$, 但
\[
 K_cp_4=cx^4-(2c+12)x^2+4,\qquad (K_cp_4)'(1)=-24,\quad
 (K_cp_4)'(-1)=24.
\]
偶函数 $K_cp_4$ 的边界值差为零, 两端导数却不为零, 所以不属于 $D(K_c)$.
按算子乘积域定义 $p_4\notin D(K_c^2)$, 与精确阈值相符.

旧版的 $N_k=(w,x^k)_t$ 任意阶取矩不成立: $x^2\notin D(K_c)$,
所以高阶内积可能未定义. 同样, $(c-D^2)^r$ 的形式多项式展开不证明
真正算子幂存在. 稠密性传输只可应用于已经属于源空间的族.
尤其从 $t=3/2$ 起, 连 $x^2$ 的 $H^t$ 矩都不能这样定义.
负阶标度虽可按上文修正, 却不消除正阶的成员障碍.
历史脚本的投影残差或形式范数增长不是幂域成员关系的证书.
$p_4$ 的成员判据也不是 $3<s<7/2$ 时整个族的稠密性证明.
本版在定理 \ref{thm:all-thresholds} 与命题 \ref{prop:h4-core} 中分别补足
全指标成员性和相容多项式逼近, 才能得出定理 \ref{thm:main}.

还须区分 $L_{\rm poly}=c-D^2$ 在全多项式空间上的代数逆与真正的 $K_c^{-1}$.
例如
\[
 L_{\rm poly}^{-1}x^2=c^{-1}x^2+2c^{-2},\qquad
 \mathcal B(L_{\rm poly}^{-1}x^2)=\binom{2/c}{-2/c}\ne0.
\]
而 $K_c^{-1}x^2$ 按定义属于 $D(K_c)$, 所以两者不相等.
本版图核心中的 $R$ 只是四维迹映射的右逆, 也不是把这个代数逆当作算子逆.

\section{仍然有效的增长引理}
\begin{lemma}[有显式系数条件的钉住解]\label{lem:growth}
设 $c_0>0$, $u_0=0$, $u_1=1$, 且
$c_0u_j=A_ju_{j-1}-B_ju_{j-2}$ ($j\ge2$).
若 $B_j\ge0$ 且 $A_j-B_j\ge c_0$ 对每个 $j\ge2$ 成立, 则
$u_j>0$ ($j\ge1$), $u_j\ge u_{j-1}$, 并且
\[
 u_j\ge\prod_{k=2}^j\frac{A_k-B_k}{c_0}.
\]
若还满足 $A_j-B_j\ge4j+c_0$, 则
$u_j\ge(4/c_0)^{j-1}j!$, 因而增长快于每个固定次数的多项式.
\end{lemma}
\begin{proof}
从 $0=u_0<u_1=1$ 出发, 归纳使用
\[
 c_0u_j=(A_j-B_j)u_{j-1}+B_j(u_{j-1}-u_{j-2})
 \ge(A_j-B_j)u_{j-1}\ge c_0u_{j-1}.
\]
逐项连乘得第一下界; 再用 $(A_j-B_j)/c_0\ge4j/c_0$ 得第二下界.
仅有前两个条件不保证快于多项式的增长: $A_j=c_0$, $B_j=0$ 时
$u_j=1$ ($j\ge1$). \qed
\end{proof}

本问题的偶、奇递推系数分别为
\[
 A_j=2j(2j-1)+\frac{cj}{j-1},\quad B_j=2j(2j-3),\qquad
 A'_j=2j(2j+1)+\frac{cj}{j-1},\quad B'_j=2j(2j-1).
\]
两者均满足 $A_j-B_j=A'_j-B'_j=4j+cj/(j-1)\ge4j+c$.
引理由直接归纳成立, 不依赖任何衰减根断言. 更确切地, 若某个递推解的
比值 $r_j=u_j/u_{j-1}$ 最终有定义并收敛到有限非零 $L$, 则
$cr_j=A_j-B_j/r_{j-1}$ 除以 $j^2$ 后给出 $0=4-4/L$, 故只能有 $L=1$.
这不声称这样的解存在或决定其增长阶. 增长引理的应用仍须先确认矩内积
有定义及其增长上界; 作为历史低阶输入的 H3 主证明只在 $H^1$ 中取矩.
本版完整窗口的主证明使用四阶图核心, 不以该增长引理替代成员性检查.

\section{数值工具的含义与限制}
\path{scripts/op10_verify_mechanism.py} 与
\path{scripts/op10_fractional_window.py} 中的有限谱和只表示谱投影的平方范数,
不自动表示原函数的有限范数. 对命题 \ref{prop:thresholds} 的发散区间,
程序应明确报告解析发散, 不以有限截断宣称通过. 在收敛区间可使用尾界:
若 $q-2t>1$, 则
\[
 \sum_{n>N}\frac{((n\pi)^2+c)^t}{(n\pi)^q}
 \le \pi^{2t-q}\left(1+\frac{c}{\pi^2(N+1)^2}\right)^{\max(t,0)}
 \frac{N^{2t-q+1}}{q-2t-1}.
\]
这是用单调因子上界和 $p$ 级数的积分尾界得到的解析估计; 普通浮点计算
该表达式仍不等于带向外舍入的区间证书.

为稳定计算第 $k$ 个正根, 写 $P_k=(k+1/2)\pi$, $\delta=P_k-\mu$,
在 $[0,\pi/2]$ 求
\[
 h_k(\delta)=\cos\delta-(P_k-\delta)\sin\delta=0.
\]
有 $h_k(0)=1$, $h_k(\pi/2)=-k\pi<0$,
$h_k'(\delta)=-(P_k-\delta)\cos\delta<0$ 于开区间, 因而严格括住唯一根.
程序检查端点符号、分支和缩放残差; 若浮点精度已不能分辨根与极点,
应显式失败, 而非把极点或固定偏移端点当根. 这些诊断、自检和有限传输
恒等式均不是独立验收或无限维稠密性的数值证明.

\begin{thebibliography}{99}
\bibitem{axioms} A. Jones, L. L. Littlejohn, A. Quintero Roba, \emph{Krein-Sobolev
	Orthogonal Polynomials II}, Axioms 14 (2025), 115.
	\url{https://doi.org/10.3390/axioms14020115}
\bibitem{lw1} L. L. Littlejohn, R. Wellman, \emph{A general left-definite theory
	for certain self-adjoint operators with applications to differential equations},
	J. Differential Equations 181 (2002), 280--339.
\bibitem{fg} C. Fischbacher, F. Gesztesy, P. Hagelstein, L. L. Littlejohn,
	\emph{Abstract left-definite theory: a model operator approach, examples,
	fractional Sobolev spaces, and interpolation theory}, arXiv:2408.01514 (2024).
	\url{https://arxiv.org/abs/2408.01514}
\bibitem{hs} 项目文档, \emph{左定空间 $H^s[-1,1]$ 中显式完备正交多项式系: 证明文档},
	2026-08-05. 算子域结论已撤回的历史文稿, 不作为本文证明前提.
\bibitem{dc} 项目文档, \emph{边界约束 Hilbert 空间中多项式稠密的一般准则},
	\path{docs/SL_denseness_criteria.tex}. 本文仅使用其原族代数包的消元论证.
\bibitem{h2} 项目文档, \emph{第二左定空间中多项式基的解析完备性: 证明文档},
	\path{docs/SL_h2_completeness_proof.tex}, 第2.1节.
\bibitem{h3} 项目文档, \emph{第三左定空间中多项式基的解析完备性},
	\path{docs/SL_h3_completeness_proof.tex}, 2026-09-20 修订的低阶主链.
\bibitem{r5} 项目文档, \emph{第二左定空间中余有限稀疏子族的精确闭包},
	\path{docs/SL_cofinite_left_definite.tex}, 第2节 R5-OP.
\bibitem{a7} 历史证明正文, \emph{Hs operator-domain vs abstract completion},
	\path{runs/three-arm-pilot-v2/pilot-v6-hs-domain/arms/c-qed/output/proof.md},
	STEP1, STEP7--STEP8. 本文只重证需要的四阶情形, 不继承该工件的审核标签.
\bibitem{r6candidate} 第六轮外部候选, \emph{原稀疏多项式族的完整非负阶范围},
	2026-09-21,
	\path{research/artifacts/proof-audit-round6-20260921/submitted/fractional_window_completion.md}.
	R6-C01 的论证来源, 非独立验收.
\end{thebibliography}

\end{document}
